This paper is the fourth in a series on the foundations of hermitian $\K$-theory using stable $\infty$-categories. In this instalment, we introduce Karoubi-localising versions of hermitian $\K$-functors and study multiplicative structures on both the original and the Karoubi-localising version.

We work with Lurie's notion of Poincaré $\infty$-categories, that is, pairs $(\C,\QF)$ where $\C$ is a stable $\infty$-category and $\QF\colon \C\op \to \Spa$ is a Poincaré structure on $\C$, that is, 
a perfect quadratic functor encoding a notion of unimodular (or Poincaré) forms in $\C$, the space of which we denote by $\Poinc(\C,\QF)$.
The fundamental invariant we consider is %
the Grothendieck-Witt functor $\GW\colon \Catp \to \Sp$, introduced in \papertwo, where the source is the $\infty$-category of Poincaré $\infty$-categories and the target the $\infty$-category of spectra.
For a ring $R$ equipped with a form parameter we in particular produced a Poincar\'e structure on the perfect derived category $\Dperf(R)$ whose Grothendieck-Witt spectrum recovers the classical Grothendieck-Witt space defined by Karoubi in \cites{karoubi-quillen,Karoubi-Le-theoreme-fondamental} as its infinite loop space, see \cite{comparison}. 

Applications of our theory given in \papertwo and \paperthree include a proof of Karoubi's periodicity conjecture, the finite generation of Grothendieck-Witt-groups for rings of integers in number fields, and explicit computations for $\ZZ$.
These computations are also a key input for forthcoming work on the stable cohomology of orthogonal and symplectic groups over the integers \cite{HLN-cohomology2,HLN-cohomologyp}. 
To obtain these applications,
we established in \papertwo what we termed the fundamental fibre sequence
\[
\K(\C,\QF)_\hC \xrightarrow{\hyp} \GW(\C,\QF) \xrightarrow{\bord} \L(\C,\QF),
\]
expressing the Grothendieck-Witt functor as an extension of the simpler algebraic $\K$- and $\L$-spectra, where the $\Ct$-action on $\K(\C)$ is induced by the duality associated to $\QF$. Since the functors $\K$ and $\L$ are Verdier-localising the same then follows for $\GW$, that is, all three take exact sequences in $\Catp$ %
to exact sequences of spectra. In fact, we showed in \papertwo that $\GW$ is the initial Verdier-localising functor under 
$\SS[\Poinc]$, %
while $\L$ is the initial such which is furthermore bordism invariant, that is, vanishes on hyperbolic Poincaré $\infty$-categories.\\


For many applications to rings and schemes, being Verdier-localising is, however, not quite sufficient. For example, on the level of perfect derived $\infty$-categories, localisations of rings yield Verdier quotients, and hence extend to exact sequences,
only up to idempotent completion. As $\K,\GW$, and $\L$ are not invariant under idempotent completion they consequently fail to satisfy Zariski descent. 
A similar problem occurs 
in the Bass-Heller-Swan and Ranicki-Shaneson decompositions concerning the $\K$- and $\L$-groups of Laurent polynomial rings, or in excision sequences for quotients by ideals.

In algebraic $\K$-theory these points were resolved with the introduction of non-connective $\K$-theory $\KK$, which is both Verdier-localising and invariant under idempotent completion, a property we call Karoubi-localising. 
Historically, there are two approaches to this extension: One, originally due to Bass and in spectral form to Thomason and Trobaugh \cites{Bass,thomason-trobaugh}, is based on the behaviour of $\K$-spectra of rings of Laurent polynomials. The other one, based on suspension rings, was suggested by Karoubi and upgraded to the spectrum level by Gersten and Wagoner in \cite{karoubi-foncteurs,gersten-spectrum, wagoner-delooping}. 
The suspension ring approach was extended by Karoubi to Grothendieck-Witt spectra, yielding what is nowadays called the Karoubi-Grothendieck-Witt functor $\KGW$,
and which indeed satisfies Zariski descent on schemes by work of Schlichting \cites{schlichting-mv,schlichting-derived}. On the other hand, Bass' construction was adapted to $\L$-spectra by Ranicki, leading to ultimately decorated $\L$-theory $\L^{\cotower{-\infty}}$, which admits well-behaved splittings for Laurent polynomial rings and %
features in %
surgery theoretic approaches to Borel's rigidity conjecture for aspherical manifolds, see 
\cite{Lueck-FJbook}. 

Unlike the case of $\KK$ and $\KGW$, which agree with $\K$ and $\GW$ respectively in positive degrees, %
the functor $\L^{\cotower{-\infty}}$ satisfies the same periodicity properties as $\L$
and 
differs from it in arbitrary degrees. In addition, while $\KK$ was recognised as universal inside Karoubi-localising functors 
on stable $\infty$-categories in~\cite{BGT}, no universal property was ever established for $\KGW$ or $\L^{\cotower{-\infty}}$.
Despite their differences, Williams conjectured in \cite{williams-quadratic} that for group rings in which 2 is invertible, $\KK$, $\KGW$, and $\L^{\cotower{-\infty}}$ should be related by an exact sequence analogous to the fundamental fibre sequence connecting the $\K$-, $\GW$- and $\L$-spectra above.

In the present paper, we extend the constructions of $\KGW$ and $\L^{\cotower{-\infty}}$ to Poincaré $\infty$-categories, characterise both by universal properties and establish Williams' conjecture in the full generality of Poincar\'e $\infty$-categories:

\begin{named}{Main theorem}
\label{theorem:pfourmaintheorem}%
For a Poincaré $\infty$-category $(\C,\QF)$ there is a natural exact sequence
\[
\KK(\C,\QF)_\hC \xrightarrow{\hyp} \KGW(\C,\QF) \xrightarrow{\bord} \L^{\cotower{-\infty}}(\C,\QF),
\]
and the functors $\KGW$ and $\L^{\cotower{-\infty}}$ enjoy the following universal properties:
\begin{enumerate}
\item
\label{item:main-theorem-i}%
they are the initial Karoubi-localising (and bordism invariant) functors $\Catp\to \Sp$ under $\SS[\Poinc]$,
respectively and
\item
\label{item:main-theorem-ii}%
they underlie the initial lax symmetric monoidal Karoubi-localising (and bordism invariant) functors $\Catp\to \Sp$, respectively.
\end{enumerate}
\end{named}


By the general results of \papertwo, 
the fibre sequence of the \reffour{theorem:pfourmaintheorem} also implies that $\L^{\langle-\infty\rangle}$ is the bordification of $\KGW$.
Moreover, given the universal properties of $\GW$ and $\L$ established in \papertwo and the general relation between bordism invariance and Karoubi-localisation, point \reffouritem{item:main-theorem-i} above shows that $\KGW$ and $\L^{\cotower{-\infty}}$ are also the initial Karoubi-localising functors under $\GW$ and $\L$, respectively. The above universal properties of $\KGW$ are thus the hermitian analogs of the universal properties of $\KK$ established in the celebrated work of Blumberg--Gepner--Tabuada \cite{BGT}. 


Our \reffour{theorem:pfourmaintheorem} can be used to infer new information on both $\KGW$ and $\L^{\cotower{-\infty}}$:
For example, one can deduce that $\KGW$ satisfies the Farrell-Jones conjecture for a given group whenever its $\KK$- and $\L^{\cotower{-\infty}}$-theoretic versions have been established for that group. This has been shown to be the case for a large class of groups by Bartels, Lück, Reich and others \cite{Lueck-FJbook}. 
At the same time, the new universal properties of $\L^{\cotower{-\infty}}$ can be used to construct a trace map from $\L^{\cotower{-\infty}}$ to the geometric fixed points of real topological cyclic homology. We again relegate the development and applications of such hermitian trace methods to forthcoming work \cite{Harpaz-Nikolaus-Shah}, building on \cites{HNS-algtrace, HNRS-algtrace}. 

We define the spectra $\L^{\cotower{-\infty}}(\C,\QF)$ by means of an inductive procedure, extending the construction for rings. Just as Bass' construction of negative $\K$-groups is based on the Bass-Heller-Swan decomposition computing the $\K$-groups of Laurent polynomial rings, the construction of $\L^{\cotower{-\infty}}$ is based on the analogous Ranicki-Shaneson decomposition. As an important ingredient in the above main theorem, we generalise such decomposition theorems to arbitrary Verdier-localising functors and arbitrary Poincar\'e $\infty$-categories, subsuming, extending and unifying these previously analogous, but not directly connected results.

Let us now describe in more detail the contents of this paper. 


\introsubsection{Karoubi-localising invariants on Poincaré \texorpdfstring{$\infty$-}{}categories}



Let us begin with our general results about Karoubi-localising approximations. Recall to that end that the functors $\GW$ and $\L$ are Verdier-localising but not Karoubi-localising.
Denoting with superscripts the full subcategories spanned by Verdier and Karoubi-localising functors, we show: 
\begin{theoremintrofour}\
\label{theorem:intro-B}%
In the solid square
\[
\begin{tikzcd}
\Fun^{\kloc}(\Catp,\Sp)\ar[r,shift right=2pt]\ar[d,"\Omega^\infty"'] & \Fun^{\vloc}(\Catp,\Sp)\ar[d,"\Omega^\infty"] \ar[l,dashed, shift right=2pt,"(-)^\mathrm{kar}"'] \\
\Fun^{\kloc}(\Catp,\Grp_{\Einf})\ar[r,shift right=2pt] & \Fun^{\vloc}(\Catp,\Grp_{\Einf}) \ar[l,dashed,shift right=2pt,"(-)^\mathrm{kar}"'] \ ,
\end{tikzcd}
\]
\begin{enumerate}
\item
\label{item:delooping}%
the left vertical functor is an equivalence, and
\item
\label{item:adjoint}%
the horizontal functors admit left adjoints $(-)^\mathrm{kar}$ making the square left adjointable.
\end{enumerate}
In particular, the canonical map $\Omega^\infty \F(-^\natural) \Rightarrow \Omega^\infty \F^{\mathrm{kar}}(-)$ is an equivalence if and only if its source is Karoubi-localising.

\end{theoremintrofour}

Theorem~\reffour{theorem:intro-B} holds equally for $\Catx$ in place of $\Catp$. In this version, \reffouritem{item:delooping} appeared previously in work of Cisinski and Khan \cite{CK}. The proof of Theorem \reffour{theorem:intro-B} also describes the horizontal adjoints and left vertical inverse explicitly: Briefly, both the successive deloopings of a Karoubi-localising $\Einf$-group-valued functor and the adjoint evaluated on a Verdier-localising spectrum-valued functor are described using (a countable version of) the Lattice-Tate sequence studied in \papertwo; it serves as a hermitian substitute of the Ind-Calkin sequence (used, for example, to produce non-connective K-theory in \cite{BGT} and continuous K-theory in \cite{Efimov}), which can itself be thought of as a categorical form of the cone-suspension ring approach of Karoubi and Villamayor. %
We also show that the left adjoint admits an alternative description using decorations, see the discussion after Theorem \reffour{theorem:intro-E} below, in our solution of Williams' conjecture.

Applying the left adjoint in Theorem \reffour{theorem:intro-B}\;\reffouritem{item:adjoint} to the spectrum-valued Verdier-localising functors $\GW$ and $\L$ we obtain Karoubi-localising approximations 
\[
\GW \Rightarrow \KGW \quad \text{ and } \quad \L \Rightarrow \KL
\]
that is, initial Karoubi-localising functors under $\GW$ and $\L$, respectively, and from the general interaction of Karoubification with bordification it follows that $\KGW^\bord \simeq \KL$. The results from \papertwo then show that $\KL$ obeys similar periodicity properties as $\L$ and yield the fundamental exact square
\[
\begin{tikzcd}
\KGW(\C,\QF) \ar[r,"\bord"] \ar[d,"\fgt"] & \KL(\C,\QF) \ar[d,"\Xi"] \\
\KK(\C,\QF)^\hC \ar[r,"\mathrm{can}"] & \KK(\C,\QF)^\tC,
\end{tikzcd}
\]
which is also the isotropy separation square of a genuine $\Ct$-equivariant non-connective real algebraic $\K$-spectrum $\KKR(\C,\QF)$.

To leverage the criterion from Theorem~\reffour{theorem:intro-B}, we consider the behaviour of $\GW$- and $\L$-spectra under Karoubi equivalences, that is, Poincar\'e functors $(f,\eta) \colon (\C,\QF) \rightarrow (\D,\QFD)$ where $f$ is fully faithful and induces an equivalence upon idempotent completion and $\eta\colon \QF \to f^*\QFD$ is an equivalence. 
As a straightforward generalisation of a result of Schlichting's in the setting of differential graded categories on which $2$ acts invertibly \cite{schlichting-derived} there are
exact sequences
\[
\GW(\C,\QF)\longrightarrow \GW(\D,\QFD) \xrightarrow{\fgt} [\K_0(\D)/\K_0(\C)]^\hC \quad \text{and} \quad \L(\C,\QF)\longrightarrow \L(\D,\QFD) \xrightarrow{\Xi} [\K_0(\D)/\K_0(\C)]^\tC.
\]
The left hand one in particular implies that the map $\GW(\C,\QF) \rightarrow \GW(\D,\QFD)$ is an isomorphism on positive homotopy groups and an injection on components, just as for $\K(\C) \rightarrow \K(\D)$; we refer to this result as the hermitian cofinality theorem. The right hand sequence is also familiar in surgery theory, entirely independent from Grothendieck-Witt theory: The associated long exact sequence is the Ranicki-Rothenberg sequence from \cite{Ranicki-lower} in the case of rings or additive categories with duality and symmetric or quadratic Poincar\'e structures.

As a consequence of the hermitian cofinality theorem, we find that $\GWspace(-^\natural)$ is Karoubi-localising, so that Theorem~\reffour{theorem:intro-B} implies that the map $\GWspace(-^\natural) \to \Omega^\infty\KGW(-)$ is an equivalence, but as a consequence of the Ranicki-Rothenberg sequences, the same is not true in $\L$-theory. To control the cofibre of the map $\L \to \KL$ (and of $\GW \rightarrow \KGW$ at the spectrum level), we finally extend the full statement of Schlichting's cofinality theorem \cite{schlichting-derived} to our setup, again removing the assumption that $2$ be invertible in the input:






\begin{propositionintrofour}
\label{proposition:intro-D}%
For any Poincaré $\infty$-category $(\C,\QF)$ the following two squares are cartesian.
\[
\begin{tikzcd}
\GW(\C,\QF) \ar[r] \ar[d,"\fgt"] & \KGW(\C,\QF) \ar[d,"\fgt"] && \L(\C,\QF) \ar[r] \ar[d,"\Xi"] & \KL(\C,\QF) \ar[d,"\Xi"]\\
\K(\C,\QF)^\hC \ar[r] & \KK(\C,\QF)^\hC &&\K(\C,\QF)^\tC \ar[r] & \KK(\C,\QF)^\tC
\end{tikzcd}
\]
\end{propositionintrofour}

In particular, if $\C$ is idempotent complete and $\KK(\C)$ happens to be connective, e.g.\ for $\C = \Dperf(R)$ with $R$ regular noetherian, we find $\GW(\C,\QF) \simeq \KGW(\C,\QF)$ and $\L(\C,\QF) \simeq \KL(\C,\QF)$. Moreover, we generally have
\[
\pi_n\cof\big(\GW(\C,\QF) \to \KGW(\C,\QF)\big)[\tfrac{1}{2}] \cong \left((\KK_n(\C)/\K_n(\C))[\tfrac 1 2]\right)^{\Ct}
\]
for all $n \in \mathbb Z$ with the right hand side being $0$ for $n>0$ and $\KK_n(\C)[\tfrac{1}{2}]^{\Ct}$ for $n<0$, while the map $\L(\C,\QF)[\frac{1}{2}] \to \KL(\C,\QF)[\frac{1}{2}]$ is an equivalence.


The above results have been used to develop Karoubi-Grothendieck-Witt-spectra in the context of schemes, where they enjoy many of the same advantages over Grothendieck-Witt spectra that non-connective $\K$-spectra have over the connective ones (and often the same is then true for $\KL$-spectra). This thread was taken up in \cite{motives} and \cite{HL}, where Zariski- and Nisnevich-descent for these spectra is established, generalising Schlichting's celebrated Mayer-Vietoris principle from \cite{schlichting-mv}, and ultimately motivic spectra (non-$\mathbb A^1$-invariant in the latter case) representing homotopical and genuine symmetric Grothendieck-Witt-theory are constructed, respectively.







\introsubsection{The decoration filtration and Williams' conjecture}

Let us now turn our attention to ultimately decorated $\L$-theory and Williams's conjecture. To explain its context, recall first the Bass-Heller-Swan %
decomposition %
\[
\KK(R[\ZZ]) \simeq \KK(R) \oplus \Sigma \KK(R) \oplus \N\KK(R)  \oplus \N\KK(R),
\]
where $R[\ZZ]$ is the group $R$-algebra of the integers and the first two summands are embedded via the assembly map $\B\ZZ \otimes \KK(R) \to \KK(R[\mathbb Z])$; the last two summands are known as the nil-terms, defined for example as the complimentary summand of $\KK(R)$ in $\KK(R[\mathbb N])$. 

Historically, the fact that in the analogous splitting of $\K(R[\mathbb Z])$ the second summand agrees with $\Sigma \K(R)$ only on connected covers, served as one of the primary motivations for introducing negative $\K$-groups. More precisely, Bass defined the spectrum $\K^{\langle -1 \rangle}(R)$ as the desuspension of this second summand and consequently $\K_{-1}(R)$ as $\pi_{-1}\K^{\langle -1 \rangle}(R)$, with the assembly map $\mathrm B\mathbb Z \otimes \K(R) \rightarrow \K(R[\mathbb Z])$ providing a comparison map $\K(R) \rightarrow \K^{\langle -1 \rangle}(R)$ that is an isomorphism in non-negative degrees. He then iterated this procedure producing spectra $\K^{\langle -n \rangle}(R)$, and $\KK(R) = \colim_{n \in \mathbb N} \K^{\langle -n \rangle}(R)$ was then the first definition of non-connective $\K$-spectra (so that, a-posteriori, $\K^{\langle -n \rangle}(R)$ coincides with $\tau_{\geq -n} \KK(R)$). 

This construction was transported to $\L$-theory to address analogous shortcomings of the Ranicki-Shaneson splitting of $\L(R[\mathbb Z])$; in fact, it simplifies slightly due to the non-occurence of nil-terms. More precisely, if $R$ is equipped with an invertible module with involution $M$ then we can equip the group ring $R[\ZZ]$ with an induced module with involution $M[\ZZ] = \oplus_{n \in \ZZ} M$ extending the involution by swapping the $n$ and $-n$ summands, and writing $\L$ for either symmetric or quadratic $\L$-spectra one puts
\[
\L^{\cotower{-1}}(R,M) \coloneq \SS^{-1} \otimes [\L(R[\ZZ],M[\ZZ])/\L(R,M)].
\]
Iterating this process again one obtains what is traditionally called the decoration sequence
\[
\L(R,M) = \L^{\langle 0\rangle}(R,M) \longrightarrow \L^{\cotower{-1}}(R,M) \longrightarrow \dots \longrightarrow \L^{\langle -n\rangle}(R,M) \longrightarrow \L^{\langle -n-1 \rangle}(R,M) \longrightarrow \dots
\]
whose colimit 
$
\L^{\cotower{-\infty}}(R,M) \coloneq \colim_{n \in \mathbb N} \L^{\cotower{-n}}(R,M)
$
is known as ultimately decorated $\L$-theory. In distinction with the case of algebraic $\K$-theory, the connecting maps in this sequence generally enjoy no particular connectivity properties. 
Despite the different constructions, in \cite{williams-quadratic} Williams predicted, in the context of groups rings over $\ZZ[\tfrac{1}{2}]$, the existence of a fibre sequence
\[
\KK(R,M)_{\hC} \to \KGW(R,M) \to \L^{\cotower{-\infty}}(R,M)
\]
involving the Karoubi-Grothendieck-Witt spectrum defined by Karoubi and Schlichting.
In the terminology of the previous subsection, his conjecture can be reformulated simply as an equivalence $\L^{\cotower{-\infty}}(R,M) \simeq \KL(R,M)$.


Our approach to proving this conjecture in general proceeds by adapting the construction above to produce an endofunctor $\mathcal T$ on $\Fun(\Catp,\Sp)$, equipped with a natural transformation from the identity, and set
\[
\F^{\langle -n \rangle} = \mathcal T^n(\mathcal F(-^\natural)) \quad \text{and} \quad \F^{\langle -\infty \rangle} = \colim_{n \in \mathbb N} \F^{\langle -n \rangle},
\]
for an arbitrary Verdier-localising $\F \colon \Catp \rightarrow \Sp$. When applied to $\L$ and $\K$ one recovers the decoration sequences described above.
But now the construction also applies to Grothendieck-Witt-spectra, where to the best of our knowledge it has not been considered before. As a simultaneous extension of Proposition~\reffour{proposition:intro-D}  and the first part of our \reffour{theorem:pfourmaintheorem} (which concern the cases $n=0$ and $n=\infty$, respectively), we then show:

\begin{theoremintrofour}
\label{theorem:intro-E}%
For any Poincaré $\infty$-category $(\C,\QF)$ and $n\geq 0$ there are exact squares
\[
\begin{tikzcd}
\GW^{\langle -n \rangle}(\C,\QF) \ar[r] \ar[d,"\fgt"] & \KGW(\C,\QF) \ar[d,"\fgt"] && \L^{\langle -n \rangle}(\C,\QF) \ar[r] \ar[d,"\Xi"] & \KL(\C,\QF) \ar[d,"\Xi"] \\
(\tau_{\geq -n} \KK(\C,\QF))^\hC \ar[r] & \KK(\C,\QF)^\hC && (\tau_{\geq -n} \KK(\C,\QF))^\tC \ar[r] & \KK(\C,\QF)^\tC
\end{tikzcd}
\]
and consequently the canonical maps $\GW^{\langle -\infty \rangle}(\C,\QF) \to \KGW(\C, \QF)$ and $\L^{\langle-\infty\rangle}(\C,\QF)\to \KL(\C,\QF)$ are equivalences.
\end{theoremintrofour}

The proof we provide in fact establishes the following generalisation of Williams conjecture: We show that the initial Karoubi-localising functor under $\F$ guaranteed by Theorem~\reffour{theorem:intro-B} 
coincides with
$\F^{\langle -\infty \rangle}$ for \emph{any} Verdier-localising $\F \colon \Catp \rightarrow \Spa$. However, in general the intermediate objects $\F^{\langle -n \rangle}$ do not seem to admit similarly simple descriptions.




\introsubsection{The circle decomposition theorem}
The workhorse result underlying Theorem \reffour{theorem:intro-E} is a unification (and far reaching generalisation) of the Bass-Heller-Swan and Ranicki-Shaneson decompositions, whose proof and discussion forms a substantial part of the present paper. To state it in full, for a functor $\G \colon \Catx \rightarrow \Sp$ and stable $\infty$-category $\D$, we set 
$\mathrm N\mathcal G(\D)$ to be the complementary summand of $\G(\C)$ in $\G(\C_{\mathrm B\mathbb N})$, where $\C_{\mathrm{B}\mathbb{N}}$ denotes the tensor of the stable $\infty$-category $\C$ with the category $\mathrm{B}\mathbb{N}$.
For example, $\Dfree(R)_{\mathrm{B}\mathbb{N}} \simeq \Dfree(R[\NN])$ while $\Dperf(R)_{\mathrm{B}\mathbb{N}}$ and $\Dperf(R[\mathbb{N}])$ still agree up to idempotent completion.
For $\G = \KK$, the nil term $\N\KK$ then coincides with the previously mentioned one when evaluated on perfect derived $\infty$-categories of rings, while for $\G=\K$ we get $\N\K = \tau_{\geq 1} \N\KK$ (which is not quite the complementary summand of $\K(R)$ in $\K(R[\NN])$, but rather the connected cover thereof). %
Replacing $\mathbb N$ by $\mathbb Z$, the operation of tensoring with $\B\ZZ$ can now also be performed on Poincaré $\infty$-categories, and one has for example an equivalence of Poincar\'e $\infty$-categories
\[
\big(\Dfree(R),\QF_M^{\mathrm{g}\lambda}\big)_{\mathrm{B}\mathbb Z}  \simeq \big(\Dfree(R[\ZZ]),\QF_{M[\ZZ]}^{\g\lambda[\ZZ]}\big)
\]
for any form parameter $\lambda=[M_{\Ct} \to Q \to M^{\Ct}]$ on $R$, where $\lambda[\ZZ]$ is the form parameter on $R[\ZZ]$ whose associated module with involution is $M[\ZZ] = \oplus_{n \in \ZZ}M$ as above (i.e.\ the involution exchanges the $n$th and $-n$th summands) and quadratic module given by $Q \oplus \big[\oplus_{n > 0}M\big]$. Note that this transforms the genuine symmetric, even, and quadratic form parameters over $R$ to genuine symmetric, even, and quadratic form parameters over $R[\ZZ]$, respectively. Similarly, one has for any spherical fibration $\xi \colon X \to \Pic(\SS)$ an equivalence of visible Poincar\'e $\infty$-categories
\[
((\Sp_{/X})^\mathrm{f},\QF^\vis_\xi)_{\mathrm B \mathbb Z} \simeq ((\Sp_{/X \times \mathrm S^1})^\mathrm{f},\QF_{\mathrm{pr}^\ast\xi}^{\vis})
\]
where $\mathrm{pr} \colon X \times \mathrm S^1 \to X$ denotes the projection.
We then generally show the following circle decomposition theorem:

\begin{theoremintrofour}
\label{theorem:intro-F}%
Let $\mathcal F \colon \Catp \rightarrow \Sp$ be Verdier-localising and $(\C,\QF)$ a Poincaré $\infty$-category. Then there is a natural equivalence
\[
\F((\C,\QF)_{\mathrm B\mathbb Z}) \simeq \F(\C,\QF) \oplus \Sigma \mathcal F(\C^\natural,\QF) \oplus \mathrm N\F^\hyp(\C).
\]
\end{theoremintrofour}


Under this decomposition the assembly map 
$
\mathrm B\mathbb Z \otimes \F(\C,\QF) \to \F((\C,\QF)_{\mathrm B\mathbb Z})
$
corresponds to the evident map
\[
\mathcal F(\C,\QF) \oplus \Sigma \mathcal F(\C,\QF) \longrightarrow \mathcal F(\C,\QF) \oplus \Sigma \mathcal F(\C^\natural,\QF) \oplus\mathrm N\F^\hyp(\C)
\]
in the first two components. %
In particular, if $\F$ is bordism invariant and $(\C,\QF)$ is idempotent complete then the assembly map
is an equivalence.


Let us now specialise the circle decomposition theorem to a few key examples.
First, for $\K \colon \Catp \rightarrow \Sp^{\mathrm{B}\Ct}$, using $\K^\hyp \simeq \Ct\otimes \K$ as functors to $\Ct$-spectra, we obtain a splitting
\[
\K((\C,\QF)_{\mathrm B\mathbb Z}) \simeq \K(\C,\QF) \oplus \Sigma \K(\C^\natural,\QF) \oplus [\Ct \otimes \mathrm{NK}(\C)],
\]
of $\Ct$-spectra. Disregarding involutions such a splitting was recently established by Saunier \cite{saunier} (starting with an arbitrary stable but idempotent complete $\D$; our proof, which is %
different from Saunier's, in fact works verbatim for all stable $\D$). If we take $\KK$ instead of $\K$ the splitting simplifies to
\[
\KK((\C,\QF)_{\mathrm B\mathbb Z}) \simeq \KK(\C,\QF) \oplus \Sigma \KK(\C,\QF) \oplus [\Ct \otimes \N\KK(\C)]
\]
as one can ignore idempotent completions. %
This recovers the classical Bass-Heller-Swan decomposition, as well as its $\mathrm A$-theoretic analogue from \cites{fundamental-spaces, fundamental-spaces2}, including an identification of the involution on the right hand side.

Second, for $\L \colon \Catp \rightarrow \Sp$, and $\KL$ using $\L^\hyp \simeq 0 \simeq \KL^\hyp$ by bordism invariance, the circle decomposition provides splittings
\[
\L((\C,\QF)_{\mathrm B\mathbb Z}) \simeq \L(\C,\QF) \oplus \Sigma \L(\C^\natural,\QF) \quad \text{ and } \quad \KL((\C,\QF)_{\mathrm B\mathbb Z}) \simeq \KL(\C,\QF) \oplus \Sigma \KL(\C,\QF).
\]
This statement generalises the Ranicki-Shaneson decomposition 
for %
quadratic and genuine symmetric $\L$-theory of ordinary rings from \cite{shaneson,Milgram-Ranicki} to arbitrary (connective in the latter case) ring spectra, and also allows to start with an arbitrary form parameter. In this regard it is perhaps of note that our approach uniformly treats all these cases. %
Let us also mention that Ranicki observed that 
there are two natural such decompositions in $\L$-theory; our general one corresponds to what he termed the algebraic decomposition. In the body of the paper we also perform a detailed analysis of the second one, the geometric decomposition, defined in terms of a certain transfer map using the visible signature of the circle. To this end we also revisit the visible signature from a new perspective in Appendix \reffour{section:appendix-visible-revisited}, which we believe to be of independent interest.



Finally, for $\GW\colon \Catp \to \Sp$, %
we obtain splittings
\[
\GW((\C,\QF)_{\mathrm B\mathbb Z}) \simeq \GW(\C,\QF) \oplus \Sigma \GW(\C^\natural,\QF) \oplus \mathrm{NK}(\C) \quad\text{and}\quad \KGW((\C,\QF)_{\mathrm B\mathbb Z}) \simeq \KGW(\C,\QF) \oplus \Sigma \KGW(\C,\QF) \oplus \N\KK(\C).
\]
that seem to be entirely new. For example, for a ring $R$ with form parameter $\lambda$ we obtain canonical equivalences
\begin{align*}
\GW\bigl(\Dfree(R[\ZZ]),\QF_{M[\ZZ]}^{\g\lambda[\ZZ]}\bigr) & \simeq \GW\bigl(\Dfree(R),\QF^{\g\lambda}_{M}\bigr) \oplus \Sigma \GW\bigl(\Dperf(R),\QF^{\g\lambda}_M\bigr) \oplus \N\K(R) \quad\text{and}\\
\KGW^{\g\lambda[\ZZ]}(R[\ZZ],M[\ZZ]) & \simeq \KGW^{\g\lambda}(R,M) \oplus \Sigma \KGW^{\g\lambda}(R,M) \oplus \N\KK(R).
\end{align*}
Furthermore, through 
the identification of the Grothendieck--Witt spectra of visible Poincar\'e $\infty$-categories with Weiss--Williams' LA-spectra from \papertwo, we in particular obtain equivalences
\begin{align*}
\LA^{\f}(X \times \mathrm S^1, \mathrm{pr}^\ast \xi) &\simeq \LA^{\f}(X, \xi) \oplus \Sigma  \LA(X, \xi) \oplus \mathrm N\mathrm A(X) \quad\text{and}\\
\mathbb{LA}(X \times \mathrm S^1, \mathrm{pr}^\ast \xi) &\simeq \mathbb{LA}(X, \xi) \oplus \Sigma  \mathbb{LA}(X, \xi) \oplus \mathrm N\mathbb A(X),
\end{align*}
where $\LA^{\f}$ is the finitely presented and $\mathbb{LA}$ the Karoubi version of the Weis-Williams $\LA$-spectra, 
which again seems to be new.


Finally, we note that Theorem~\reffour{theorem:intro-F} may also be applied to further invariants such as real (topological) Hochschild and real (topological) cyclic homology, once their domain of definition is extended to Poincar\'e $\infty$-categories in \cite{Harpaz-Nikolaus-Shah}.

\begin{comment}
\introsubsection{Karoubification of Verdier-localising functors}

The results explained  so far do not only establish the exact sequence in the \reffour{theorem:pfourmaintheorem} along with the first universal property, but they also show that two, a priori very different constructions, of $\KGW$ and $\KL$ agree:  The first one takes the colimit of the decoration tower whose $n$-th  entry is $\F^{\langle-n\rangle}$, which can be formed a priori for an arbitrary functor $\F\colon \Catp\to \Spa$. This construction is a categorical and hermitian analogue of Bass's original iterative construction of negative K-theory. 

The  other construction is implicit in the inverse map of Theorem \reffour{theorem:intro-B} and is the colimit of what we call the \defi{Tate-tower}, which is a tower of the form
\[
\F(\C, \QF) \to \Omega \F((\C, \QF)\otimes (\T, \QFtwo)) \to \Omega^2 \F((\C, \QF)\otimes (\T, \QFtwo) \otimes (\T, \QFtwo)) \to \dots
\]
where $(\T, \QFtwo)=\Tate^{\omega_1}(\Spaf, (\QF^\uni)\qshift 1)$ is a countable version of the Tate Poincar\'e category encountered in \papertwo; here, $\F=\GW, \L$ but it can in principle also be applied to an arbitrary functor $\F\colon \Catp\to \Spa$. In a first approximation, this construction may be thought of as a categorical and hermitian analogue of the suspension ring construction of negative K-theory. 

As an addendum to this part of our \reffour{theorem:pfourmaintheorem}, we show that both towers produce universal Karoubi-localising approximations even in the context of arbitrary (in particular, not necessarily finitary) Verdier-localising functors:

\begin{theoremintrofour}
\label{theorem:intro-F'}%
Let $\F\colon \Catp \to \Sp$ be any Verdier-localising functor. Then the colimits of both the Tate-tower and the decoration tower are universal Karoubi-localising approximations for $\F$. In particular, the inclusion 
\[
\Fun^\kloc(\Catp, \Spa) \to \Fun^\vloc(\Catp, \Spa)
\]
from Karoubi-localising into Verdier-localising functors admits a left adjoint.
\end{theoremintrofour}

As explained, the  cases $\F=\GW, \L$ of this result also follow from the results established before; this line  of reasoning uses certain connectivity estimates which are ultimately consequences of the hermitian cofinality theorem, and so does not generalize to the general case. For instance, it is not even clear, a priori, that any of the two towers converges to a Verdier-localising functor; the point is that the individual stages of the tower are additive and Karoubi invariant, but can fail to have localisation sequences in any range of degrees.

We mention that Theorem \reffour{theorem:intro-F'} fails with ``Verdier-localising'' replaced  by ``additive'', and that we do not know whether the inclusion from Karoubi-localising into all additive functors (in the absence of any finiteness assumptions) admits a left adjoint.
\end{comment}







\introsubsection{Multiplicative structures on Grothendieck-Witt- and L-spectra}
As the final set of results which we establish in this paper, we turn to multiplicative structures. Our main result in this direction, subsuming in particular the final assertion of the \reffour{theorem:pfourmaintheorem}, is the following: 
\begin{theoremintrofour}
\label{theorem:intro-G}%
The cube
\[
\begin{tikzcd}
[column sep=0.5em, row sep=0.7em]
  & \KGW\ar[rr]\ar[dd] && \KL\ar[dd]\\
\GW\ar[ru]\ar[rr, crossing over]\ar[dd] && \L\ar[ur]\\
  & \KK^\hC\ar[rr] && \KK^\tC\\
\K^\hC\ar[rr] \ar[ru] && \K^\tC\ar[uu,<-, crossing over] \ar[ru]\\
\end{tikzcd}
\]
refines to a cube of lax symmetric monoidal functors $\Catp \rightarrow \Sp$, in a way that makes the functors $\GW$ and $\KGW$ into the initial Verdier- and Karoubi-localising lax symmetric monoidal functors $\Catp \to \Sp$, respectively, and similarly $\L$ and $\KL$ into the initial ones that are in addition bordism invariant. 
\end{theoremintrofour}


In particular, if $(\C,\QF)$ is a (commutative) algebra in $\Catp$, that is, if $\C$ is a stably (symmetric) monoidal $\infty$-category and $\QF \colon \C\op \rightarrow \Sp$ is a lax (symmetric) monoidal Poincaré structure, then for any of the functors $\mathcal F$ in the cube above, the spectrum $\F(\C,\QF)$ acquires the structure of a (commutative) ring spectrum. Likewise, if $(\D,\Phi)$ is a $(\C,\QF)$-module in $\Catp$, then $\F(\D,\Phi)$ acquires the structure of an $\F(\C,\QF)$-module spectrum. 
For example, this applies when $\C = \Dperf(R)$ for a commutative ring $R$ and $\QF$ is one of $\QF^\sym_R$ or $\QF_R^\mathrm{gs}$, recovering and extending results of Laures and McClure from \cite{laures} on multiplications for $\L$-spectra (though we do not 
compare our structures to theirs), and when $\D=\Dperf(A)$ for $A$ an $R$-algebra and $\Phi$ one of $\QF^\qdr_A$ or $\QF^\mathrm{gq}_A$. We spell out precise conditions on form parameters encoding that one is a module over another in Appendix \reffour{section:appendix-form-parameters}. The front face of the cube above was also very recently refined to a lax symmetric monoidal one in \cite{HSV26} (and here our universal property immediately implies compatibility with their structure); our results on the back face, in particular for $\LL = \L^{\langle -\infty \rangle}$ seem to be new. 


We go on to show that the association $(\C,\QF) \mapsto [(\C,\QF^\qdr) \to (\C,\QF)]$ refines to a lax symmetric monoidal functor $\Catp \to \Fun([1],\Catp)$, where the target is equipped with Day convolution with respect to the minimum on $[1]$.
It then follows that the functor $\NL \colon \Catp \rightarrow \Sp$, given by the common value (as a consequence of the fundamental fibre sequence and Proposition \reffour{proposition:intro-D}) of 
\[
(\C,\QF) \longmapsto \cof(\F(\C,\QF^\qdr) \rightarrow \F(\C,\QF))
\]
for any of $\F = \GW, \KGW, \L$ or $\KL$ referred to as normal or hyperquadratic $\L$-theory by Ranicki, inherits a lax symmetric monoidal structure. Combining such arguments with the fact that $\NL_0(\Dperf(\mathbb Z),\QF^\mathrm{gs}) \cong \mathbb Z/8$, one obtains for example:

\begin{corollaryintrofour}
\label{corollary:intro-H}%
For any discrete ring $R$ and invertible module with involution $M$ over $R$, the kernel and cokernel of each of the polarisation maps
\[
\GW_i^\qdr(R,M) \longrightarrow \GW_i^\mathrm{gq}(R,M) \longrightarrow \GW_i^\mathrm{gs}(R,M) \longrightarrow \GW_i^\mathrm{s}(R,M)
\]
 and their composites is $8$-torsion for every $i \in \ZZ$, and similarly for $\KGW$, $\L$, or $\KL$ in place of $\GW$.
\end{corollaryintrofour}

This statement really does not extend to arbitrary form parameters or (connective) ring spectra, as Weiss' and Williams' explicit calculations for the sphere spectrum from \cite{WWIII} and our calculation for the genuine Frobenius parameter from \paperthree show.

Our proof of Theorem~\reffour{theorem:intro-G} relies on a general analysis of the Day convolution $\infty$-operads $\Fun(\C,\D)^\Day$ and their localisation properties, to which we devote Appendix \reffour{section:appendix-localisations}. The colours of this operad are the functors $\C \to \D$ and its commutative algebra objects are the lax symmetric monoidal such functors. Our arguments are then based on the following observations:
\begin{enumerate}
\item
The Day convolution $\infty$-operad $\Fun(\Catp,\Spa)^{\Day}$, though not a symmetric monoidal $\infty$-category due to size issues, still admits a unit object, given explicitly by $\SS[\Poinc]$.
\item
Each of the functors $\GW,\KGW,\L,\KL$ is initial under $\SS[\Poinc]$ with respect to a certain property: being additive, Karoubi-localising, or a combination of one of these with bordism invariance. Each of these properties can be phrased as sending a certain collection $\mathcal A$ of cone diagrams to limit diagrams.
\item
When the collection $\mathcal A$ is suitably compatible with the tensor product on $\Catp$, such universal approximations of the unit in $\Fun(\Catp,\Spa)^{\Day}$ become themselves the units of the full suboperad spanned by the functors which send all diagrams in $\mathcal{A}$ to limit diagrams. In particular, these universal approximations refine, in a suitably unique manner, to initial commutative algebra objects, hence initial lax symmetric monoidal functors of the given type. %
\item When $\mathcal A$ consists of finite diagrams and is furthermore stable under forming tensors and cotensors with arbitrary stable $\infty$-categories, these commutative algebra objects are even idempotent and therefore yield smashing localisations onto the category of functors $\Catp \to \Spa$ taking all diagrams in $\mathcal A$ to limit diagrams. 
\end{enumerate}

The last point applies to the collections of diagrams defining additivity and bordism invariance, giving:

 \begin{comment}

To treat all cases uniformally we cast our arguments in a general setting as follows.
Consider two symmetric monoidal $\infty$-categories $\C$ and $\D$, where $\D$ admits colimits distributing over tensor products (e.g., $\D=\Sps,\Spa$). 
While Day convolution does not make $\Fun(\C,\D)$ into a symmetric monoidal $\infty$-category
in this generality,
it does give rise to an $\infty$-operad $\Fun(\C,\D)^\Day$ admitting a unit object $\one_\Day = [c \mapsto (\one_\D)_{\Hom_\C(\one_\C,c)}]$, %
which consequently refines to the initial lax symmetric monoidal functor $\C \to \D$. Now suppose given a collection $\cR$ of cone diagrams $K^\triangleleft \rightarrow \C$ (with possibly varying $K \in \Cat$), and denote by $\Fun^\cR(\C,\D)$ the full subcategory of $\Fun(\C,\D)$ on $\cR$-local functors, that is, those sending the diagrams in $\cR$ to limit diagrams in $\D$.
We then formulate convenient criteria under which the initial $\cR$-local approximation $\one_\cR$ to $\one_\Day$, should it exist, refines, uniquely under $\one_\Day$, to the initial lax symmetric monoidal $\cR$-local functor. This is the case, for example, if the collection $\cR$ is stable under tensoring with arbitrary objects from $\C$ and $\D$ is %
closed symmetric monoidal. Let us note that if $\C$ is small and $\D$ is presentable the Day convolution $\infty$-operad is a presentably symmetric monoidal $\infty$-category and $\Fun^{\cR}(\C,\D)$ is a left Bousfield localisation of $\Fun(\C,\D)$, in which case the above statement follows by fairly standard arguments concerning multiplicative left Bousfield localisations. To treat the general case we have dedicated Appendix \reffour{section:appendix-localisations} %
to the analysis of multiplicative localisations in the setting of $\infty$-operads, %
which we believe to be useful in many other contexts beyond the present one.

Returning to the case $\C=\Catp$ and $\D=\Spa$, 
the above mentioned results are sufficient to establish the top face of the cube in Theorem \reffour{theorem:intro-G}, 
as well as the multiplicativity properties of $\K$ and $\KK$. To obtain the entire cube one then relies on the multiplicativity properties of the adjunction $U\colon \Catp \adj \Catx\cocolon \Hyp$. %
As the mechanism behind these results is fairly general, further applications such as trace invariant functors on laced $\infty$-categories, or $\mathbb{A}^1$-invariant functors on schemes, are easily envisioned. 

In certain favourable situations the $\cR$-local unit turns out to not just be lax symmetric monoidal but idempotent with respect to Day convolution, in a manner such that Day convolving with it yields a smashing localisation of $\Fun(\C,\D)$ into $\Fun^{\cR}(\C,\D)$. Of the functors mentioned in Theorem~\reffour{theorem:intro-G}, this applies to $\K,\GW$ and $\L$, but not to $\KK,\KGW$ and $\KL$. The precise sufficient conditions needed to insure the smashing property are summarized in the following result:

\begin{propositionintrofour}
\label{proposition:intro-I}%
Let $\C$ and $\D$ be closed symmetric monoidal categories with $\D$ cocomplete and $\cR$ a collection of cone diagrams in $\C$. Assume the following hypotheses hold:
\begin{enumerate}
\item
The collection $\cR$ is stable under tensoring and cotensoring with objects in $\C$. 
\item
The class of $\cR$-local functors is closed under colimits and tensoring with objects of $\D$.
\item
The unit $\one_{\Day}$ admits an initial $\cR$-local approximation $\one_{\Day} \Rightarrow \one_{\cR}$ such that $\one_{\cR}$ is a small functor (that is, left Kan extended from a small subcategory of $\C$).
\end{enumerate}
Then $\one_{\cR}$ is idempotent with respect to Day convolution and convolving with it yields a smashing localisation of $\infty$-operads
\[
\Fun(\C,\D)^{\Day} \longrightarrow \Fun^\cR(\C,\D)^{\Day}, \quad G \longmapsto \one_{\cR} \otimes_\Day G .
\]
In particular, in this case $\Fun^\cR(\C,\D)^{\Day}$ identifies with the $\infty$-operad of $\one_{\cR}$-modules in $\Fun(\C,\D)^{\Day}$.
\end{propositionintrofour}


Proposition~\reffour{proposition:intro-I} appears to be new even in the case where $\C$ is small. For example, it implies that it $\C$ is a closed semi-additively symmeric monoidal small $\infty$-category $\C$ and $\D$ is a semi-additive presentably symmetric monoidal $\infty$-category then the full subcategory $\Fun^{\oplus}(\C,\D)$ of direct sum preserving functors is a smashing localisation of $\Fun(\C,\D)$. Similarly,   
if $\C$ is a closed stably symmetric monoidal small $\infty$-category and $\D$ is a %
closed and cocomplete stable symmetric monoidal $\infty$-category then the full subcategory $\Exc^n_*(\C,\D)$ of reduced $n$-excisive functors $\C \to \D$ is a smashing localisation of $\Fun(\C,\D)$. In the case of $\C=\Spaf$ and $\D=\Spa$, the latter was already observed in~\cite{ABHS}. To put matters in perspective, we note that even $\Exc^1_*(\C,\Spa)$ is typically not a smashing localisation when $\C$ is not stable, see Example~\reffour{example:cant-remove} in Appendix \reffour{section:appendix-localisations}.

Now while neither the collection of Verdier- nor Karoubi-sequences is closed under the cotensoring of $\Catp$ over itself, the collection of split Poincaré-Verdier sequences 
is so, and the same holds for the collection of bordism equivalences in $\Catp$. 
We therefore conclude:
\end{comment}

\begin{propositionintrofour}
\label{proposition:intro-J}%
Both functors $\GW,\L \colon \Catp \to \Sp$ are idempotent with respect to Day convolution and convolving with them yields a smashing localisation onto the full subcategory of additive functors $\Catp \to \Sp$ in the case of $\GW$ and additive bordism invariant functors $\Catp \rightarrow \Sp$ in the case of $\L$. %
\end{propositionintrofour}


A completely analogous statement holds for $\K \colon \Catx \rightarrow \Sp$
and it follows that all functors from $\Catp$ or $\Catx$ to spectra admit additive (and bordism invariant) approximations. By contrast, the collection defining Karoubi-localising functors is not closed under cotensors and indeed we do not know whether any of the functors $\KK, \KGW$ or $\KL$ is idempotent. Following a suggestion of Ramzi, we do, however, exhibit non-Karoubi-localising modules over all three of these functors in the body of the paper, which rules out that Karoubi-localising functors form a smashing localisation of all functors $\Catp \rightarrow \Spa$. We consequently do not know whether all functors to spectra admit Karoubi-localising approximations (with or without bordism invariance).

There are many variants of the above that can be studied, for example by considering the functors $\K$, $\GW$, $\L$ and their Karoubi-localisations also after postcomposition with the forgetful functors $\Sp \to \Grp_{\Einf} \to \Sps$. For instance, we show that the functor $\Lspace\colon\Catp \to \Grp_{\Einf}$ is idempotent and localises onto bordism invariant and direct sum preserving functors (in particular bordism invariant additive functors to $\Grp_{\Einf}$ also do not form a smashing localisation inside all functors). In upcoming work, Saunier shows that $\Kspace\colon \Catx \to \Sps$ and $\GWspace\colon \Catp \to \Grp_{\Einf},\Sps$ are also idempotent and identifies the associated localised subcategories.



Let us finally mention one application outside the realm of hermitian K-theory, as a further advertisement of our localisation machinery: If $\C$ is a closed stably symmetric monoidal small $\infty$-category and $\D$ is a %
closed and cocomplete stable symmetric monoidal $\infty$-category then our criteria apply to show that the full subcategory $\Exc^n_*(\C,\D)$ of reduced $n$-excisive functors $\C \to \D$ is a smashing localisation of $\Fun(\C,\D)$. The case of $\C=\Spaf$ and $\D=\Spa$ was recently obtained in~\cite{ABHS} and to put matters into perspective, note that even $\Exc^1_*(\C,\Spa)$ is typically not a smashing localisation when $\C$ is not stable. %





\introsubsection{Organisation of the paper}
In \S\reffour{section:localising}, we discuss Karoubi-localising functors and prove %
the second part of Theorem~\reffour{theorem:intro-B} in Theorem ~\reffour{proposition:localising-are-stable}, and Proposition~\reffour{proposition:intro-D} as Theorem~\reffour{theorem:Karoubi-pullback-squares}, respectively. \S\reffour{section:shaneson} is then devoted to the circle decomposition theorem, Theorem~\reffour{theorem:intro-F}, and the subsequent \S\reffour{section:karoubification} etablishes Theorem~\reffour{theorem:intro-E}, and thus Williams' conjecture as Propositions~\reffour{proposition:decorated-GW} and~\reffour{proposition:decorated-L} along with the remaining part of Theorem~\reffour{theorem:intro-B} in Corollary~\reffour{corollary:karoubification}.
Finally, \S\reffour{section:day} discusses multiplicative structures on all Grothendieck-Witt and L-functors constructed in this paper. We prove Theorem~\reffour{theorem:intro-G} as Corollary~\reffour{corollary:kgw-kl-monoidal}, Corollary~\reffour{corollary:intro-H} as Corollary~\reffour{corollary:normalL8tors}, and Proposition~\reffour{proposition:intro-J} in Corollaries~\reffour{corollary:GW-universal-multiplicative} and~\reffour{corollary:lmonoidal}. 

Finally, the paper has four appendices. In \S\reffour{section:appendix-localisations}, we discuss multiplicative and smashing localisations of $\infty$-operads which lay the foundations of our results in \S\reffour{section:day}. %
In \S\reffour{section:presentability} we show that $\Catp$ and $\Cath$ are compactly generated, a technical result we use a few times along the way, in \S\reffour{section:appendix-visible-revisited} we revisit the visible symmetric signature of a Poincar\'e complex to facilitate the discussion of the two possible splittings in the circle decomposition theorem, and in \S\reffour{section:appendix-form-parameters} we unwind in detail our mulitplicitativity results into the more classical language of form parameters. %
